\chapter{Grammatical Gender, to Mastery}
\label{ch:gender-mastery}

\emph{Practice units: \texttt{units/ES-P0-U16} through \texttt{U21}.}

This is the chapter every \emph{(masc.)}/\emph{(fem.)} tag since Chapter~1
has been building toward.

\section{Articles}

\begin{vocabtable}
el / la & --- & the (masc./fem. singular) \\
los / las & --- & the (masc./fem. plural) \\
un / una & --- & a/an (masc./fem.) \\
\end{vocabtable}

\begin{etymology}
\emph{el/la} $\leftarrow$ Latin \emph{ille/illa} --- the same root as
\emph{\'{e}l/ella} (Chapter~1). \emph{un/una} $\leftarrow$ Latin
\emph{unus/una}, the same root as the number \emph{uno}.
\end{etymology}

\begin{grammarlens}
English articles (\emph{the}, \emph{a}) never change form. Spanish
articles must match the noun's gender and number: \emph{el libro, la
casa, los libros, las casas}. This is the first hard evidence that gender
ripples outward to everything describing a noun, starting with its
article.
\end{grammarlens}

\section{The Pattern --- and Its Real Exceptions}

The default: nouns ending in \textbf{-o} are usually masculine, nouns
ending in \textbf{-a} are usually feminine. ``Usually'' is doing real work:

\begin{vocabtable}
el d\'{i}a & masc. & day (ends in -a, masculine anyway) \\
la mano & fem. & hand (ends in -o, feminine anyway) \\
el problema & masc. & problem (Greek \emph{-ma} noun, masculine) \\
el agua & fem. & water (feminine noun, masculine-sounding article) \\
\end{vocabtable}

\begin{etymology}
\emph{d\'{i}a} $\leftarrow$ Latin \emph{dies}, masculine in the Latin form
Spanish inherited from --- the gender came whole from Latin, not from the
later \emph{-o/-a} generalization. \emph{mano} $\leftarrow$ Latin
\emph{manus}, feminine in its Latin declension class despite the \emph{-o}
ending it ended up with in Spanish. \emph{problema, programa, sistema,
tema} are Greek-origin \emph{-ma} nouns that kept a masculine gender
despite the \emph{-a} spelling. \emph{agua} is grammatically feminine but
takes \emph{el} purely to avoid two stressed ``a'' sounds colliding ---
adjectives describing it still agree feminine (\emph{el agua fr\'{i}a}).
\end{etymology}

\begin{grammarlens}
Grammatical gender is a noun-classification system \textbf{inherited from
Latin}, not a rule invented from scratch. The \emph{-o/-a} pattern is a
later simplification that fits most nouns, but wherever a word's Latin (or
Greek) origin had a different gender, that original gender usually
survived --- which is exactly why gender has to be memorized per word
rather than computed from spelling.
\end{grammarlens}

\section{Adjective Agreement}

\begin{vocabtable}
alto / alta & --- & tall \\
peque\~{n}o / peque\~{n}a & --- & small \\
grande & --- & big (same form, both genders) \\
\end{vocabtable}

\begin{grammarlens}
Spanish adjectives agree in gender and number with their noun: \emph{el
libro alto, la casa alta, los libros altos, las casas altas}. English
adjectives never change form (``tall book,'' ``tall house''). A handful of
adjectives like \emph{grande}, ending in \emph{-e} or a consonant, don't
change by gender at all --- only by number.
\end{grammarlens}

\section{Colors}

\begin{vocabtable}
rojo / roja & --- & red \\
azul & --- & blue \\
verde & --- & green \\
amarillo / amarilla & --- & yellow \\
negro / negra & --- & black \\
\end{vocabtable}

\begin{etymology}
\emph{azul} $\leftarrow$ Arabic \emph{l\={a}zaward} --- the same root
behind English \emph{azure}, both languages borrowing independently from
Arabic, another entry alongside \emph{cero} in Spanish's Arabic thread.
\emph{verde} $\leftarrow$ Latin \emph{viridis} (English \emph{verdant}).
\emph{negro} $\leftarrow$ Latin \emph{niger} (English \emph{denigrate}).
\end{etymology}

\section{Family Vocabulary}

\begin{vocabtable}
la madre & fem. & mother \\
el padre & masc. & father \\
el hermano / la hermana & --- & brother / sister \\
el hijo / la hija & --- & son / daughter \\
\end{vocabtable}

\begin{etymology}
\emph{madre}/\emph{padre} $\leftarrow$ Latin \emph{mater}/\emph{pater}
(English \emph{maternal, paternal, patron}). \emph{hijo/hija} $\leftarrow$
Latin \emph{filius/filia} (English \emph{filial, affiliate}).
\end{etymology}

\begin{grammarlens}
Family words are the one vocabulary set where grammatical gender and
real-world sex line up directly --- \emph{padre} is masculine because
fathers are male, not because of the \emph{-o/-a} pattern (it ends in
\emph{-e}, outside that pattern entirely). Worth noticing precisely
because it's the exception: most Spanish nouns have gender with zero
connection to real-world sex.
\end{grammarlens}

Chapter 3 closes with a capstone describing a family, recombining
articles, the gender pattern, adjective agreement, colors, and family
vocabulary in one description (\texttt{units/ES-P0-U21}).
